\documentclass[]{article}

%opening
\title{HPC PCB Testing apparatus and methods}
\author{Nyle Systems Control}

\begin{document}

\maketitle

\begin{abstract}
The HPC PCB is a product used in Nyle heat pump units, manufactured by JLPCB. The boards require testing to ensure quality. The purpose of this document is to describe how this process may be performed.
The Testing Apparatus will provide all of the connections necessary to test the HPC, including power to the board. The HPC Testing Apparatus is and electro mechanical device intended only to perform quality controls against a cross section of HPC PCBs created. 
\end{abstract}

\section*{Its Automated!}
The purpose of the HPC-TA  is to perform the 100 to 1000 tests that may guaranty the reasonable quality, of the HPC-TA with minimal human intervention. Never the less, it is a human operated device. Human operators must remain aware of possible fault conditions. A short circuit, for example could cause fire, or a hot smell. Any irregular behaviour from the HPC-TA should result in it being turned off. 

\section*{TP: Test Points}
The HPC PCB Contains 221 TP, 1-236, enumerated as TP1, TP2, ... TP236. The Test points allow for component testing, and live circuit testing. For the 
\subsection*{Types of Tests}
Power: GND, 24V, and Voltage Regulator Tests. 
Component tests: This may include continuity to gnd, or specified test points.
Powered Component tests: This may involve tests of components where power is required to conduct the tests, such as but not limited to switching relays, blinking LEDs, or communicating with specific devices beyond basic connectivity tests. 

\section*{Methods}

\section*{}
\section*{}
\section*{}

\end{document}
